% SGA 3, Expose VI_A, Loop 1 English reconstruction.
% Scope: paragraphs 5.5 and 5.6, through Proposition 5.6.2 and its
% proof in full; Section 6 is excluded.
% Authority: Polo--Gille Exp6A-13oct24.pdf, component-local pages 33--34
% (combined-reader pages 336--337; printed pages 325--326).
% OCR and the Reinhold English Markdown are comparison witnesses only.
% Loop 2 replaces the two exact 180-dpi source crops with native tikz-cd
% reconstructions controlled against the PDF, local page 34:
%   quotient square: x=605,y=710,w=300,h=155 on a 1488x2105 render;
%   local-flatness square: x=575,y=1095,w=350,h=150 on that render.

\medskip
\phantomsection\label{par:sga3-VIA-5.5}%
\noindent\textbf{5.5.---}
设 \(G\) 是在 Artin 局部环 \(A\) 上局部有限型且平坦的群.
由 \SGARefLink{cor:sga3-VIB-2.3}{2.3} 已知,
单位连通分支 \(G^0\) 是 \(G\) 的一个正规开子群概形,
因而也在 \(A\) 上平坦. 于是由 \hyperref[thm:sga3-VIA-3.2]{3.2} 和
\hyperref[par:sga3-VIA-5.2]{5.2} 可知,
\[
  G^0\backslash G
\]
是一个在 \(A\) 上平坦的 \(A\)-群概形.
此外, 由于 \(G\) 的每个连通分支 \(G^\alpha\) 对 \(G^0\) 定义的等价关系饱和,
商 \(G^0\backslash G\) 是各 \(G^0\backslash G^\alpha\) 的不交并
(参见 \hyperref[par:sga3-VIA-4.3]{4.3}).
特别地, \(G^0\backslash G\) 的单位连通分支恰好是
\[
  G^0\backslash G^0\simeq\operatorname{Spec}A,
\]
因而态射
\[
  G^0\backslash G\longrightarrow\operatorname{Spec}A
\]
在单位元处是局部同构. 因此, 由第~VI\(_B\)~篇
\SGARefLink{prop:sga3-VIB-1.3}{1.3},
\(G^0\backslash G\) 在 \(\operatorname{Spec}A\) 上平展.%
\footnote{\label{note:sga3-VIA-55}编者注: 对原文前面的讨论作了扩充,
并增添了\hyperref[prop:sga3-VIA-5.5.1]{命题~5.5.1}, 以明确陈述这一结果.}
因此得到如下命题 (断言~\SGARefLink{prop:sga3-VIA-5.5.1:item:ii}{(ii)}
见 \textup{[DG70]}, \S~II.5, 1.7--1.10).

\medskip
\phantomsection\label{prop:sga3-VIA-5.5.1}%
\noindent\textbf{命题 5.5.1.---}
\emph{设 \(A\) 是 Artin 局部环,
\(G\) 是局部有限型且平坦的 \(A\)-群.
\begin{enumerate}
\item[\textup{(i)}] \(G^0\backslash G\) 是平展 \(A\)-群.
\item[\textup{\SGAInlineRefTarget{prop:sga3-VIA-5.5.1:item:ii}(ii)}]
因此, 若 \(A=k\) 为代数闭域,
则 \(G^0\backslash G\) 是一个在 \(G\) 的连通分支集上单传递作用的常 \(k\)-群.
所以, 若 \(G\) 为代数群, 则 \(G^0\backslash G\) 有限.
\end{enumerate}}

\medskip
\phantomsection\label{par:sga3-VIA-5.6}%
\noindent\textbf{5.6.---}
现设 \(k\) 为完美域, \(G\) 为局部有限型 \(k\)-群.
我们在 \SGARefLink{sga3:VIA:0.2}{0.2} 中已看到,
此时 \(G_{\mathrm{red}}\) 为 \(G\) 的子群概形.
此外, 在 \(G_{\mathrm{red}}\) 对 \(G\) 的左作用下,
\(G\) 的单位元所在的等价类就是 \(G\) 的整个底空间.
因此\hyperref[thm:sga3-VIA-3.2]{定理~3.2} 给出下述结果.

\medskip
\phantomsection\label{prop:sga3-VIA-5.6.1}%
\noindent\textbf{命题 5.6.1.---}
\emph{设 \(k\) 为完美域, \(G\) 为局部有限型 \(k\)-群.
则 \(k\)-概形
\[
  G_{\mathrm{red}}\backslash G
\]
是剩余域为 \(k\) 的一个有限局部 \(k\)-代数的谱.}

\noindent
\footnote{\label{note:sga3-VIA-56}编者注: 编号 5.6.1 及下面的证明为增补内容.}
事实上, 由 \SGARefLink{thm:sga3-VIA-3.2}{3.2},
商 \(G_{\mathrm{red}}\backslash G\) 只有一个点, 剩余域为 \(k\),
且它是有限型 \(k\)-概形.
因此, 它是某个有限维局部 \(k\)-代数的谱 (参见 EGA~I, 6.4.4).

\medskip
\phantomsection\label{prop:sga3-VIA-5.6.2}%
\noindent\textbf{命题 5.6.2.---}
\emph{设 \(u:F\to G\) 是完美域 \(k\) 上局部有限型群之间的态射.
下述断言等价:
\begin{enumerate}
\item[\textup{\SGAInlineRefTarget{prop:sga3-VIA-5.6.2:item:i}(i)}] \(u\) 平坦.
\item[\textup{\SGAInlineRefTarget{prop:sga3-VIA-5.6.2:item:ii}(ii)}] 态射
\[
  u^0:F^0\longrightarrow G^0
\]
占优, 且 \(u\) 诱导的态射
\[
  v:F_{\mathrm{red}}\backslash F
    \longrightarrow G_{\mathrm{red}}\backslash G
\]
平坦.
\end{enumerate}}

\noindent
\footnote{\label{note:sga3-VIA-57}编者注: 我们扩充了
\textup{\SGARefLink{prop:sga3-VIA-5.6.2:item:ii}{(ii)}} \(\Rightarrow\)
\textup{\SGARefLink{prop:sga3-VIA-5.6.2:item:i}{(i)}} 的证明,
并简化了下图.}
事实上, 考虑如下交换图:
\begin{center}
  \phantomsection\label{diag:sga3-VIA-5.6.2-quotients}%
  \begin{tikzcd}[row sep=large,column sep=large]
    F
      \arrow[r,"p"]
      \arrow[d,"u"'] &
    F_{\mathrm{red}}\backslash F
      \arrow[d,"v"] \\
    G
      \arrow[r,"q"'] &
    G_{\mathrm{red}}\backslash G
  \end{tikzcd}
\end{center}
其中 \(p\) 和 \(q\) 为典范投影.
由 \SGARefLink{concl:sga3-VIA-3.2-iv}{3.2(iv)},
\(p\) 和 \(q\) 忠实平坦. 因此, 若 \(u\) 平坦, 则
\[
  q\circ u=v\circ p
\]
平坦, 因而 \(v\) 也平坦.

反过来, 假设 \(v\) 平坦且 \(u^0\) 占优.
由于 \(u^0\) 拟紧 (由 \SGARefLink{prop:sga3-VIA-2.4}{2.4},
概形 \(F^0\) 在 \(k\) 上有限型, 因而为 Noether 概形),
它将 \(F^0\) 的泛点 \(\xi\) 映到 \(G^0\) 的泛点 \(\eta\).
设 \(R\) 为满足
\[
  G_{\mathrm{red}}\backslash G=\operatorname{Spec}R,
\]
的有限局部 \(k\)-代数, 并设 \(\mathfrak m\) 为其极大理想.
有局部环的局部同态
\[
  R\longrightarrow\mathcal O_{G,\eta}
   \longrightarrow\mathcal O_{F,\xi}.
\]
注意下述方块为笛卡尔方块:
\begin{center}
  \phantomsection\label{diag:sga3-VIA-5.6.2-local-flatness}%
  \begin{tikzcd}[row sep=large,column sep=large]
    G_{\mathrm{red}}
      \arrow[r]
      \arrow[d] &
    G
      \arrow[d,"q"] \\
    \operatorname{Spec}(R/\mathfrak m)
      \arrow[r] &
    \operatorname{Spec}(R)
  \end{tikzcd}
\end{center}
因此
\[
  \mathcal O_{G,\eta}/\mathfrak m\mathcal O_{G,\eta}
    \simeq \mathcal O_{G_{\mathrm{red}},\eta}
    =\kappa(\eta),
\]
从而
\[
  \mathcal O_{F,\xi}/\mathfrak m\mathcal O_{F,\xi}
\]
在 \(\mathcal O_{G,\eta}/\mathfrak m\mathcal O_{G,\eta}\) 上平坦.

另一方面, 由于 \(q\) 和 \(v\circ p\) 平坦,
\(G\) 和 \(F\) 在 \(R\) 上平坦.
因此, 平坦性的局部判据 (参见 EGA~IV\(_3\), 11.3.10.2) 表明,
\(\mathcal O_{F,\xi}\) 在 \(\mathcal O_{G,\eta}\) 上平坦;
换言之, \(u\) 在 \(\xi\) 处平坦.
由 \hyperref[lem:sga3-VIA-2.5.3]{2.5.3} 可知 \(u\) 平坦.
% Independent prose review: source and translation compared paragraph by paragraph through this component; statements, hypotheses, proofs, notes and diagram captions covered; bibliography reserved for its assigned reviewer; source math and reference keys unchanged.
